\documentclass[11pt,a4paper]{article}

\usepackage[margin=25mm]{geometry}
\usepackage[T1]{fontenc}
\usepackage{lmodern}
\usepackage{microtype}
\usepackage{amsmath,amssymb,amsthm,mathtools}
\usepackage{booktabs,array,longtable}
\usepackage[numbers,sort&compress]{natbib}
\usepackage[hidelinks]{hyperref}
\usepackage{enumitem}
\setlist{nosep}

\newtheorem{theorem}{Theorem}
\newtheorem{proposition}[theorem]{Proposition}
\newtheorem{corollary}[theorem]{Corollary}
\newtheorem{lemma}[theorem]{Lemma}
\theoremstyle{definition}
\newtheorem{definition}[theorem]{Definition}
\newtheorem{remark}[theorem]{Remark}

\newcommand{\ZZ}{\mathbb Z}
\providecommand{\T}{}
\renewcommand{\T}{\mathcal T}
\providecommand{\A}{}
\renewcommand{\A}{\mathcal A}
\newcommand{\Ozero}{\mathcal O_0}
\newcommand{\Ocomp}{\mathcal O_{\rm comp}}
\newcommand{\Pcomp}{\mathcal P_{\rm comp}}
\newcommand{\HTp}{\mathrm{HT}_{+}}
\newcommand{\HTm}{\mathrm{HT}_{-}}

\title{An Exactly Solvable Family of Routed Rotation Systems\\
with Observation-Dependent Future Quotients}
\author{Maks Katsubo\\Independent researcher}
\date{Publication-grade consolidation v0.6 --- 6 October 2026}

\begin{document}
\maketitle

\begin{abstract}
We define an infinite family of finite routed rotation systems
\[
\T_{n;\delta,c_{AB},c_{BC}},\qquad n=4q\ge 8,
\]
obtained by a cyclic deformation of a geometry-derived PIFI routing template.
The full system has \(7n-1\) vertices, \(16n\) undirected edges and \(32n\) directed edge states, but its non-radial dynamics admits an exact \(2n\)-state return reduction.
The induced macro shift is
\[
s=2(\delta+c_{AB}+c_{BC})\pmod n,
\]
and, after reduction modulo \(q\), the return dynamics is a canonical two-layer arithmetic system \(\A(q,s)\).
Writing \(g=\gcd(q,s)\) and \(m=q/g\), we obtain closed formulas for the coarse cycle census, primitive future periods, future-equivalence quotient size \(Q\), and resolving depth \(k_{\min}\).

The original PIFI slice \((\delta,c_{AB},c_{BC})=(1,1,1)\) has \(s=6\); in that slice
\[
Q=52\iff n\in\{8,12,24\},
\]
and the original \(n=12\) system has \(Q=52\) and \(k_{\min}=13\).

We then separate graph microstructure from observable future behavior.
The AB- and BC-family restricted graphs carry component invariants \(h_{AB}\) and \(h_{BC}\), while the coarse future quotient depends only on \(s\).
Refining the output alphabet by exact family-component labels yields an exact decorated-orbit formula
\[
Q_{\rm comp}
=
8+\sum_{[W]\in\Pcomp}|W|.
\]
The refined quotient depends on component placement, not merely on component counts.
That placement admits the normal form
\[
(n,s,d_{AB},d_{BC},\rho),\qquad \rho=c_{BC}\bmod d_{AB},
\]
and, modulo an affine index gauge, the remaining datum is the diagonal unit orbit
\(
[(s,\rho)]_{U(n)}.
\)
Finally, this orbit is classified prime by prime by truncated \(p\)-adic valuations and a relative unit ratio on the common nonzero prime-power layer.

The general machinery---rotation systems, follower/future equivalence, cyclic words, affine actions and CRT---is standard.
The candidate contribution is the exact solvability and end-to-end arithmetic classification of the explicit routed family defined here.
\end{abstract}

\medskip
\noindent\textbf{Keywords:}
rotation systems; symbolic dynamics; future equivalence; cyclic words; affine actions; Chinese remainder theorem; computational combinatorics.

\section{Introduction}

Finite graph dynamics often becomes much smaller after one passes from the full state space to an induced return system or to an output-equivalence quotient.
The general principles behind such reductions are classical: combinatorial maps encode cyclic port order on darts, transition systems specify local continuation rules, and follower or right-language equivalence identifies states with the same future output sequence
\cite{Moore1956,LindMarcus2021,Johansen2011}.
Likewise, cyclic lifts and arithmetic orbit decompositions naturally produce gcd and order phenomena
\cite{GrossTucker1977,HubardMochanMontero2023}.

This paper studies one explicit family for which all of those layers can be calculated exactly.
The family originates from a fixed PIFI geometric template, but the main statements below are purely finite and combinatorial.
The geometric motivation singles out the original point
\[
n=12,\qquad
(\delta,c_{AB},c_{BC})=(1,1,1),
\]
where a previously observed future quotient has \(52\) states.
Rather than treating \(52\) as an isolated numerical phenomenon, we place that point inside an infinite arithmetic family.

The logical chain is
\[
(\delta,c_{AB},c_{BC})
\longrightarrow
s
\longrightarrow
\A(q,s)
\longrightarrow
(Q,k_{\min},\text{cycle census}),
\]
followed by
\[
\text{component labels}
\longrightarrow
\text{decorated macro orbits}
\longrightarrow
\text{affine placement}
\longrightarrow
\text{prime-power CRT data}.
\]

\subsection{Claim boundary}

No new theory of combinatorial maps, automata minimization, sofic covers, voltage graphs, primitive necklaces, affine permutations, finite abelian groups, or the Chinese remainder theorem is claimed.
Those frameworks are established
\cite{LindMarcus2021,Johansen2011,GrossTucker1977,HubardMochanMontero2023,BorsWang2022,DuttaPrasad2011,Cao2011}.
The novelty candidate is restricted to the exactly solvable family defined below, its explicit reductions, and their combined classification.

A targeted prior-art audit did not locate the complete theorem chain presented here.
This is evidence for a family-specific novelty candidate, not a proof of publication priority.

\section{The routed rotation-system family}\label{sec:def}

Fix
\[
n=4q,\qquad q\ge2,
\]
and read every index modulo \(n\).
The parameters satisfy
\[
\delta\not\equiv0,\qquad
c_{AB}\not\equiv0,\qquad
c_{BC}\not\equiv0,\frac n2
\pmod n.
\]

\begin{definition}[Underlying labelled graph]
Before seam identification the vertices are
\[
\{O\}\cup
\{A_i,B_i,C_i,X_i,Y_i,Z_i,W_i:i\in\ZZ_n\}.
\]
Impose
\[
C_0\sim C_{2q}=:D_x,\qquad
C_q\sim C_{3q}=:D_y.
\]
The five edge families are
\begin{align*}
OA:\;& OX_i,\ X_iA_i,\\
AA:\;& A_iY_i,\ X_iY_i,\ A_iY_{i-\delta},\ X_iY_{i-\delta},\\
AB:\;& A_iZ_i,\ A_{i+c_{AB}}Z_i,\ Z_iB_i,\ Z_iB_{i+c_{AB}},\\
AC:\;& A_iB_i,\ B_iC_i,\\
BC:\;& B_iW_i,\ B_{i+c_{BC}}W_i,\ W_iC_i,\ W_iC_{i+c_{BC}}.
\end{align*}
\end{definition}

\begin{proposition}[Graph counts]\label{prop:counts}
Every admissible system has
\[
\boxed{V=7n-1},\qquad
\boxed{E=16n},\qquad
\boxed{|D|=32n},
\]
where \(D\) is the dart set, and degree histogram
\[
\boxed{
(n-4)d_3+4n\,d_4+(2n+2)d_6+d_n.
}
\]
\end{proposition}

\begin{proof}
Before the seam quotient there are \(7n+1\) vertices; the two independent identifications reduce the count by two.
The five edge families contribute \(2n+4n+4n+2n+4n=16n\) edges.
The degree count follows directly from the incidence table.
\end{proof}

\subsection{Rotation data and routing}

The cyclic port order is part of the definition:
\[
\begin{array}{c|l}
\text{vertex} & \text{cyclic neighbor order}\\
\hline
O & X_0,X_1,\ldots,X_{n-1}\\
X_i & Y_{i-\delta},O,Y_i,A_i\\
Y_i & X_i,X_{i+\delta},A_{i+\delta},A_i\\
Z_i & A_i,A_{i+c_{AB}},B_{i+c_{AB}},B_i\\
W_i & B_i,B_{i+c_{BC}},C_{i+c_{BC}},C_i\\
A_i & Y_{i-\delta},Z_{i-c_{AB}},B_i,Z_i,Y_i,X_i\\
B_i & Z_{i-c_{AB}},A_i,Z_i,W_i,C_i,W_{i-c_{BC}}.
\end{array}
\]
For non-cardinal \(C_i\),
\[
C_i:\quad B_i,W_i,W_{i-c_{BC}}.
\]
At the two seam vertices:
\begin{align*}
D_x:\;&B_0,W_0,W_{-c_{BC}},B_{2q},W_{2q},W_{2q-c_{BC}},\\
D_y:\;&B_q,W_q,W_{q-c_{BC}},B_{3q},W_{3q},W_{3q-c_{BC}}.
\end{align*}

At every even-degree vertex, an incoming dart leaves by the opposite port.
At every degree-three vertex, choose the next port in one cyclic direction for \(\HTp\) and the previous port for \(\HTm\).
This defines a permutation
\[
P:D\to D.
\]

\begin{definition}[Coarse observation]
For a dart \(d:u\to v\) belonging to family \(F\), define
\[
\Ozero(d)=(F,\deg u,\deg v).
\]
Two darts are future-equivalent when
\[
d\sim_\infty d'
\quad\Longleftrightarrow\quad
\Ozero(P^t d)=\Ozero(P^t d')
\ \text{for every }t\ge0.
\]
Let \(Q\) be the number of equivalence classes.
The resolving depth \(k_{\min}\) is the least \(k\) for which equality of the first \(k\) observed symbols determines the complete future class.
\end{definition}

\section{Exact macro reduction}\label{sec:macro}

The non-radial dynamics is captured by
\[
\boxed{L_j=W_j\to B_j},
\qquad
\boxed{R_j=W_j\to B_{j+c_{BC}}}.
\]
Let
\[
S=\{0,q,2q,3q\}\subset\ZZ_n.
\]

\begin{theorem}[Macro-shift theorem]\label{thm:macro}
Define
\[
\boxed{s=2(\delta+c_{AB}+c_{BC})\pmod n.}
\]
For \(\HTp\),
\[
L_j\longmapsto
\begin{cases}
R_{j+n/2-(s-c_{BC})},&j-(s-c_{BC})\in S,\\
L_{j-s},&j-(s-c_{BC})\notin S,
\end{cases}
\]
and
\[
R_j\longmapsto
\begin{cases}
L_{j+n/2+s-c_{BC}},&j+s\in S,\\
R_{j+n/2+s},&j+s\notin S.
\end{cases}
\]
The output-block lengths are \(12\) on regular \(L\)-frames, \(20\) on regular \(R\)-frames, and \(12\) on seam-switch frames. We denote these three coarse block types by \(C3\), \(C20\), and \(C6\), respectively; their explicit output words are listed in Appendix~\ref{app:frames}.
\end{theorem}

\begin{proof}
Trace an \(L_j\) or \(R_j\) entry through the fixed local rotation table.
The AA, AB and BC offsets contribute twice to the complete frame displacement, giving \(2(\delta+c_{AB}+c_{BC})\).
The only layer switches occur when the trace hits one of the identified cardinal \(C\)-positions.
The intermediate local degrees and edge families are fixed, hence so are the frame lengths.
\end{proof}

Set
\[
x=j-(s-c_{BC})\quad\text{on }L,
\qquad
y=j+s\quad\text{on }R.
\]
Modulo \(q\), the four cardinal indices become \(0\).

\begin{definition}[Canonical arithmetic return system]
Let \(\A(q,s)\) be the permutation of
\(
\{L_x,R_y:x,y\in\ZZ_q\}
\)
defined by
\[
L_0\mapsto R_s,\qquad
L_x\mapsto L_{x-s}\ (x\ne0),
\]
\[
R_0\mapsto L_{-s},\qquad
R_y\mapsto R_{y+s}\ (y\ne0).
\]
\end{definition}

\begin{corollary}[Canonical reduction]\label{cor:Aqs}
The non-radial routed dynamics of
\(
\T_{n;\delta,c_{AB},c_{BC}}
\)
reduces exactly to \(\A(q,s)\).
Therefore all coarse future invariants depend on the three wiring parameters only through \(s\).
\end{corollary}

\section{Coarse arithmetic classification}\label{sec:coarse}

Set
\[
\boxed{g=\gcd(q,s),\qquad m=\frac qg.}
\]
The mixed core word is
\[
\boxed{C6\,C20^{\,m-1}\,C6\,C3^{\,m-1}.}
\]

\begin{theorem}[Primitive future classes]\label{thm:future}
The primitive coarse future periods are
\[
\begin{cases}
8,\ 32m-8,&g=1,\\[1mm]
8,\ 12,\ 12,\ 20,&m=1,\\[1mm]
8,\ 12,\ 20,\ 32m-8,&g>1,\ m>1.
\end{cases}
\]
Consequently
\[
\boxed{
Q=
\begin{cases}
32q,&g=1,\\[1mm]
52,&m=1,\\[1mm]
32(m+1),&g>1,\ m>1.
\end{cases}}
\]
\end{theorem}

\begin{proof}
Away from the seam the two layers are translations by \(-s\) and \(+s\) modulo \(q\), with orbit length \(m=q/g\).
A mixed orbit meets the seam twice and traverses one \(C6\), then \(m-1\) \(C20\)-frames, another \(C6\), and \(m-1\) \(C3\)-frames.
Its primitive output length is
\[
12+20(m-1)+12+12(m-1)=32m-8.
\]
The three cases follow from whether pure translation orbits avoiding the seam exist.
Summing distinct primitive periods gives \(Q\).
\end{proof}

\begin{theorem}[Resolving depth]\label{thm:kmin}
The exact resolving depth is
\[
\boxed{
k_{\min}=
\begin{cases}
13,&m=1,\\
25,&g=1,\ m=2,\\
20m-29,&g=1,\ m\ge3,\\
20m-9,&g>1,\ m>1.
\end{cases}}
\]
\end{theorem}

This is the maximal common-prefix length among distinct cyclic phases of the primitive periodic futures, plus one.

For the full routed cycle census let
\[
d_L=\gcd(n,s),\qquad
d_R=\gcd(n,n/2+s).
\]
Then
\[
N_{C3}=d_L-\frac{d_L}{\gcd(d_L,q)},
\qquad
L_{C3}=12\frac n{d_L},
\]
\[
N_{C20}=d_R-\frac{d_R}{\gcd(d_R,q)},
\qquad
L_{C20}=20\frac n{d_R},
\]
and
\[
N_M=
\begin{cases}
4,&m\text{ odd},\\
2,&m\text{ even},
\end{cases}
\qquad
L_M=
\begin{cases}
32m-8,&m\text{ odd},\\
2(32m-8),&m\text{ even}.
\end{cases}
\]
Four radial \(U\)-cycles of length \(8\) are always present.

\section{The original PIFI slice}\label{sec:pifi}

The geometry-derived wiring has
\[
\boxed{(\delta,c_{AB},c_{BC})=(1,1,1),}
\]
so
\[
\boxed{s=6.}
\]

\begin{corollary}[The 52-state exceptional set]\label{cor:52}
Along this wiring slice,
\[
\boxed{
Q=52
\iff
q\mid6
\iff
n\in\{8,12,24\}.
}
\]
For the original \(n=12\) system,
\[
\boxed{Q=52,\qquad k_{\min}=13,}
\]
with primitive periods
\[
\boxed{8,\ 12,\ 12,\ 20.}
\]
\end{corollary}

\section{Information loss under the coarse observation}\label{sec:loss}

\begin{proposition}[AB component invariant]\label{prop:hAB}
The AB-only subgraph has
\[
\boxed{h_{AB}=\gcd(n,c_{AB})}
\]
connected components.
\end{proposition}

\begin{proof}
AB incidence changes the cyclic index only by multiples of \(c_{AB}\), hence connected components are cosets of the subgroup generated by \(c_{AB}\) in \(\ZZ_n\).
\end{proof}

Set
\[
d_{BC}=\gcd(n,c_{BC}).
\]

\begin{proposition}[BC component invariant]\label{prop:hBC}
After the two cardinal \(C\)-identifications, the BC-only subgraph has
\[
\boxed{
h_{BC}=
\begin{cases}
d_{BC},&d_{BC}\mid n/2,\\
d_{BC}-2,&d_{BC}\nmid n/2.
\end{cases}}
\]
\end{proposition}

\begin{proposition}[Explicit many-to-one coarse quotient]\label{prop:manytoone}
At \(n=12\), the four triples
\[
(1,1,1),\quad
(1,3,5),\quad
(1,5,3),\quad
(1,4,4)
\]
all have
\[
s=6,\qquad Q=52,\qquad k_{\min}=13,
\]
but respectively
\[
(h_{AB},h_{BC})
=
(1,1),\ (3,1),\ (1,3),\ (4,2).
\]
\end{proposition}

Thus the coarse output alphabet forgets genuine labelled-graph microstructure.

\section{Refined observation and component-decorated orbits}\label{sec:refined}

Define
\[
d_{AB}=\gcd(n,c_{AB}),
\qquad
\alpha(i)=i\bmod d_{AB}.
\]
For BC, start with residues modulo
\(
d_{BC}
\)
and impose the two quotient relations induced by
\[
0\sim n/2,\qquad q\sim3q.
\]
Let
\[
\kappa:\ZZ_{d_{BC}}\to\mathcal C_{BC}
\]
be the resulting component map.

\begin{definition}[Component-aware observation]
For an AB dart belonging to the \(Z_i\)-cell, refine \(\Ozero\) by \(\alpha(i)\).
For a BC dart belonging to the \(W_i\)-cell, refine \(\Ozero\) by \(\kappa(i)\).
All other symbols remain unchanged.
Call the refined observation \(\Ocomp\) and its quotient size \(Q_{\rm comp}\).
\end{definition}

At the original PIFI point
\(
h_{AB}=h_{BC}=1
\),
hence the added labels are trivial.

\begin{corollary}[PIFI robustness]
For the original \(n=12\) PIFI system,
\[
\boxed{Q_{\rm comp}=Q=52,}
\]
with
\[
k_{\min}=13,\qquad
8,12,12,20.
\]
\end{corollary}

For the four \(n=12,s=6\) examples from Proposition~\ref{prop:manytoone}, the refined sizes are
\[
52,\ 96,\ 84,\ 184.
\]

Use the macro gauge
\[
x=j-(s-c_{BC})\quad\text{on }L,
\qquad
y=j+s\quad\text{on }R,
\]
and put
\[
\boxed{\rho=c_{BC}\bmod d_{AB}.}
\]
Then the AB component labels inside a frame are
\[
L_x:\quad x+s-\rho,\ x+\rho\pmod{d_{AB}},
\]
\[
R_y:\quad y-s+\rho,\ y-\rho\pmod{d_{AB}}.
\]
The BC labels are
\[
L/C3:\quad \kappa(x),\kappa(x),
\]
\[
L/C6:\quad \kappa(x),\kappa(x+n/2),
\]
\[
R/C20\text{ or }C6:\quad \kappa(y),\kappa(y+n/2).
\]

Let \(\beta(z)\) be the corresponding decorated output block at macro state \(z\).

\begin{theorem}[Component-orbit formula]\label{thm:componentorbit}
Let \(\Omega\) be the cycles of the exact \(2n\)-state macro permutation.
For
\[
C=(z_0,\ldots,z_{\ell-1})\in\Omega
\]
form
\[
W_C=\beta(z_0)\beta(z_1)\cdots\beta(z_{\ell-1}).
\]
Let \(\Pcomp\) be the distinct cyclic classes of the primitive roots
\(
\operatorname{prim}(W_C)
\).
Then
\[
\boxed{
Q_{\rm comp}
=
8+\sum_{[W]\in\Pcomp}|W|.
}
\]
\end{theorem}

\begin{proof}
The radial sector contributes one primitive class of length \(8\).
Every non-radial orbit returns to a macro entry and emits exactly the concatenated word \(W_C\).
Two periodic futures are equivalent precisely when their primitive decorated words agree up to cyclic phase.
\end{proof}

Examples include
\[
52=8+12+12+20,
\]
\[
84=8+12+12+12+20+20,
\]
\[
96=8+12+12+20+20+24,
\]
\[
184=8+4\cdot12+4\cdot20+2\cdot24.
\]

\begin{proposition}[Counts alone are insufficient]\label{prop:192312}
There is no formula
\[
Q_{\rm comp}=F(n,s,h_{AB},h_{BC})
\]
for this refined observation.
At \(n=24\), the systems
\[
(1,7,6)
\quad\text{and}\quad
(1,5,8)
\]
both have
\[
s=4,\qquad(h_{AB},h_{BC})=(1,6),
\]
but
\[
\boxed{Q_{\rm comp}=192}
\qquad\text{and}\qquad
\boxed{Q_{\rm comp}=312}.
\]
\end{proposition}

\section{Affine placement normal form}\label{sec:placement}

\begin{theorem}[Raw placement normal form]\label{thm:rawplacement}
Up to component-label renaming, the component-labelled macro system is determined by
\[
\boxed{
\Pi_0=(n,s,d_{AB},d_{BC},\rho).
}
\]
\end{theorem}

The cardinal seam set is preserved by affine index relabellings
\[
j\mapsto uj+t\pmod n,
\]
where \(U(n)=(\ZZ/n\ZZ)^\times\), \(u\in U(n)\), and \(t\in\{0,q,2q,3q\}\).
After component-label renaming, \(t\) is immaterial, while
\[
(s,\rho)\mapsto
(us\bmod n,u\rho\bmod d_{AB}).
\]

\begin{theorem}[Affine placement class]\label{thm:affine}
The affine-gauge placement class is
\[
\boxed{
\Theta=
\left(
n,d_{AB},d_{BC},[(s,\rho)]_{U(n)}
\right).
}
\]
Equivalently, with
\[
\epsilon_{BC}=
\begin{cases}
1,&d_{BC}\mid n/2,\\
0,&d_{BC}\nmid n/2,
\end{cases}
\]
one may write
\[
\Theta=
\left(
n,h_{AB},h_{BC},\epsilon_{BC},[(s,\rho)]_{U(n)}
\right).
\]
\end{theorem}

Pure gcd data do not classify the pair \((s,\rho)\): the coupled residue orbit is essential.

\section{Prime-power CRT classification}\label{sec:crt}

Let
\[
d=d_{AB},
\qquad
n=\prod_p p^{e_p},
\qquad
d=\prod_p p^{f_p}.
\]
For each prime \(p\mid n\), define
\[
\boxed{a_p=\min(v_p(s),e_p),}
\qquad
\boxed{b_p=\min(v_p(\rho),f_p).}
\]
If
\[
a_p<e_p,\qquad b_p<f_p,
\]
set
\[
\boxed{m_p=\min(e_p-a_p,f_p-b_p)}
\]
and
\[
\boxed{
r_p=
\left(\frac{\rho}{p^{b_p}}\right)
\left(\frac{s}{p^{a_p}}\right)^{-1}
\pmod{p^{m_p}}.
}
\]
If either local coordinate vanishes, set \(m_p=0\) and omit \(r_p\).

\begin{theorem}[Prime-power classification]\label{thm:crt}
Two pairs
\[
(s,\rho),(s',\rho')\in\ZZ_n\times\ZZ_d
\]
lie in the same diagonal \(U(n)\)-orbit if and only if, for every \(p\mid n\),
\[
\boxed{
(a_p,b_p,m_p,r_p)
=
(a'_p,b'_p,m'_p,r'_p).
}
\]
No exceptional \(p=2\) case is required.
\end{theorem}

\begin{proof}
A unit preserves the truncated valuations and cancels in the relative unit ratio, proving necessity.
For sufficiency, after removing the fixed powers of \(p\), a common multiplier must satisfy
\[
u\equiv s'_0s_0^{-1}\pmod{p^{e_p-a_p}},
\qquad
u\equiv \rho'_0\rho_0^{-1}\pmod{p^{f_p-b_p}}.
\]
These congruences are compatible exactly modulo
\(
p^{m_p}
\),
and equality of \(r_p\) is precisely the compatibility condition.
A local unit therefore exists for every prime, and CRT combines them into one \(u\in U(n)\).
\end{proof}

Hence the final arithmetic placement signature is
\[
\boxed{
\Theta_{\rm CRT}
=
\left(
n,d_{AB},d_{BC},
\{(a_p,b_p,m_p,r_p)\}_{p\mid n}
\right).
}
\]

This is an elementary specialization of standard \(p\)-primary orbit ideas for finite abelian groups and finite chain rings
\cite{DuttaPrasad2011,Cao2011}; no new general orbit theory is claimed.

\section{Final logical boundary and controls}\label{sec:boundary}

The CRT signature is necessary and sufficient for the diagonal affine placement orbit:
\[
\boxed{
\Theta_{\rm CRT}=\Theta'_{\rm CRT}
\iff
[(s,\rho)]_{U(n)}
=
[(s',\rho')]_{U(n)}.
}
\]

The component-colour future language can be coarser.
Final controls found distinct CRT placement classes with identical component-colour future languages after component-label renaming.
Therefore the valid implication is
\[
\boxed{
\text{same CRT placement class}
\Longrightarrow
\text{same component-colour future language},
}
\]
but not generally the converse.

\begin{table}[h]
\centering
\caption{Final control summary.}
\begin{tabular}{lrr}
\toprule
control & cases & mismatches\\
\midrule
full graph HT+ vs.\ macro theorem & 152 & 0\\
coarse full graph, both chiralities & 304 & 0\\
refined full graph vs.\ component-orbit predictor & 152 & 0\\
refined HT+/HT- quotient invariants & 304 & 0\\
generic CRT pair theorem, \(2\le n\le64\) & 146965 & 0\\
larger-modulus CRT samples & 4335 & 0\\
extended adversarial placement cases, \(n\le240\) & 144 & 0\\
\bottomrule
\end{tabular}
\end{table}

These computations guard against implementation errors and certify finite instances.
They do not establish publication priority.

\section{Discussion}

The main result is not the isolated value \(52\).
The stronger statement is that a \(32n\)-dart routed family collapses to a two-layer arithmetic return system whose coarse future quotient is governed by one shift \(s\), while controlled observation refinement exposes additional component-placement information.

The original PIFI point is robust in one precise sense:
its AB and BC restricted graphs are connected, so component refinement adds no information and leaves
\[
Q=52.
\]
Deformed systems with the same coarse quotient can expand to
\[
84,\ 96,\ 184,\ 192,\ 248,\ 312,\ldots
\]
under the refined observation.

The resulting hierarchy is
\[
\boxed{
\text{routed graph}
\to
\text{macro return}
\to
\text{coarse future quotient}
\to
\text{refined decorated orbits}
\to
\text{affine placement}
\to
\text{CRT data}.
}
\]

Rotation systems, future/follower equivalence, cyclic-word reduction, affine actions and \(p\)-primary decomposition are established tools
\cite{LindMarcus2021,Johansen2011,GrossTucker1977,BorsWang2022,DuttaPrasad2011,Cao2011}.
The candidate contribution is their exact synthesis for the explicit family \(\T_{n;\delta,c_{AB},c_{BC}}\).

\section{Reproducibility}

All theorem-level claims are accompanied by exact Python verifiers and machine-readable JSON certificates.
The publication package freezes the chain:

\begin{enumerate}
\item split-annular graph and macro theorem;
\item refined terminal injectivity;
\item component-orbit predictor;
\item affine placement classification;
\item CRT prime-power classifier;
\item final necessity/sufficiency control;
\item specialist novelty/claim-boundary audit.
\end{enumerate}

A public snapshot should export only the manuscript, these frozen verifiers and recursive local dependencies, the corresponding result certificates, and the claim-boundary documents.

\appendix

\section{Explicit coarse frame words}\label{app:frames}

For completeness, write \(F|a|b\) for an observed dart in edge family \(F\)
with tail degree \(a\) and head degree \(b\).
The regular \(L\)-frame, seam-switch frame, and regular \(R\)-frame have
the following coarse words.

The \(12\)-symbol regular \(L\) word is
\[
\begin{aligned}
C3={}&(
AB|6|4,\ AB|4|6,\ AA|6|4,\ AA|4|4,\ AA|4|4,\ AA|4|6,\\
&AB|6|4,\ AB|4|6,\ BC|6|4,\ BC|4|3,\ BC|3|4,\ BC|4|6).
\end{aligned}
\]

The \(12\)-symbol seam-switch word is
\[
\begin{aligned}
C6={}&(
AB|6|4,\ AB|4|6,\ AA|6|4,\ AA|4|4,\ AA|4|4,\ AA|4|6,\\
&AB|6|4,\ AB|4|6,\ BC|6|4,\ BC|4|6,\ BC|6|4,\ BC|4|6).
\end{aligned}
\]

The \(20\)-symbol regular \(R\) word is
\[
\begin{aligned}
C20={}&(
AB|6|4,\ AB|4|6,\ AA|6|4,\ AA|4|4,\ AA|4|4,\ AA|4|6,\\
&AB|6|4,\ AB|4|6,\ BC|6|4,\ BC|4|3,\ AC|3|6,\ AC|6|6,\\
&OA|6|4,\ OA|4|n,\ OA|n|4,\ OA|4|6,\ AC|6|6,\ AC|6|3,\\
&BC|3|4,\ BC|4|6).
\end{aligned}
\]

The radial primitive word is
\[
\begin{aligned}
U={}&(
OA|6|4,\ OA|4|n,\ OA|n|4,\ OA|4|6,\\
&AC|6|6,\ AC|6|6,\ AC|6|6,\ AC|6|6).
\end{aligned}
\]

These explicit words make the period sums in
Theorem~\ref{thm:future} and Theorem~\ref{thm:componentorbit}
directly checkable from the stated macro orbit structure.

\section{Evidence map}


\footnotesize
\begin{longtable}{p{0.31\textwidth}p{0.31\textwidth}p{0.31\textwidth}}
\toprule
claim & result certificate & verifier / document\\
\midrule
split-annular macro theorem &
\path{docs/results/pifi_3d_tn_split_annular_offset_v0_1.json} &
\path{tools/verify_pifi_3d_tn_split_annular_offset_v0_1.py}\\
refined observation / injectivity &
\path{docs/results/pifi_3d_refined_terminal_injectivity_v0_1.json} &
\path{tools/verify_pifi_3d_refined_terminal_injectivity_v0_1.py}\\
component-orbit formula &
\path{docs/results/pifi_3d_refined_quotient_component_orbit_v0_1.json} &
\path{tools/verify_pifi_3d_refined_quotient_component_orbit_v0_1.py}\\
affine placement normal form &
\path{docs/results/pifi_3d_component_placement_affine_gauge_v0_1.json} &
\path{tools/verify_pifi_3d_component_placement_affine_gauge_v0_1.py}\\
CRT prime-power theorem &
\path{docs/results/pifi_3d_crt_prime_power_placement_v0_1.json} &
\path{tools/verify_pifi_3d_crt_prime_power_placement_v0_1.py}\\
final control &
\path{docs/results/pifi_3d_final_classification_control_v0_1.json} &
\path{tools/verify_pifi_3d_final_classification_control_v0_1.py}\\
novelty boundary &
\path{docs/results/pifi_3d_specialist_novelty_literature_audit_v0_1.json} &
\path{docs/research/pifi_2d_square_ca/PIFI_3D_SPECIALIST_NOVELTY_LITERATURE_AUDIT_v0_1.md}\\
\bottomrule
\end{longtable}
\normalsize

\bibliographystyle{plainnat}
\bibliography{references}

\end{document}
